\RequirePackage[thinlines]{easybmat} %-- muss aufgrund eines Bugs vor etex und tikz geladen werden
\documentclass[a4paper, twoside, headsepline, index=totoc,toc=listof, fontsize=10pt, cleardoublepage=empty, headinclude, DIV=13, BCOR=13mm, titlepage,]{scrartcl}
\usepackage{scrtime} % KOMA, Uhrzeit ermoeglicht
\usepackage{scrpage2} % wie fancyhdr, nur optimiert auf KOMA-Skript, leicht andere Syntax
\usepackage{etoolbox}
\usepackage{letltxmacro}

% Compiler abhaengige IFs
\usepackage{ifxetex,ifluatex}
\newif\ifxetexorluatex
\ifxetex
  \xetexorluatextrue
\else
  \ifluatex
    \xetexorluatextrue
  \else
    \xetexorluatexfalse
  \fi
\fi

%--Farbdefinitionen
%-- muss vor tikz geladen werden
\usepackage[usenames, table, x11names]{xcolor}
\definecolor{dark_gray}{gray}{0.45}
\definecolor{light_gray}{gray}{0.6}

%--Zum Zeichnen
%-- muss vor polyglossia bzw. babel geladen werden
\usepackage{tikz}
\usepackage{tikz-cd}
\usetikzlibrary{external}
\tikzset{>=latex}
\usetikzlibrary{shapes,arrows,intersections}
\usetikzlibrary{calc,3d}
\usetikzlibrary{decorations.pathreplacing,decorations.markings}
\usetikzlibrary{angles}

%-- Konfiguration von tikzexternalize
\tikzexternalize[prefix=tikz/, up to date check=diff]

\ifxetexorluatex
\tikzset{external/system call={xelatex \tikzexternalcheckshellescape %-- verwende LuaLaTeX, wegen dynamischer Speicherallokation
    -halt-on-error -interaction=batchmode -jobname "\image" "\texsource"}}
\else
\tikzset{external/system call={pdflatex \tikzexternalcheckshellescape 
    -halt-on-error -interaction=batchmode -jobname "\image" "\texsource"}}
\fi

%-- tikzexternalize fuer tikzcd deaktivieren, da inkompatibel
\AtBeginEnvironment{tikzcd}{\tikzexternaldisable}
\AtEndEnvironment{tikzcd}{\tikzexternalenable}

%-- um Inkompatibilitaeten von quotes und polyglossia bzw. babel zu vermeiden
\tikzset{
  every picture/.append style={
    execute at begin picture={\shorthandoff{"}},
    execute at end picture={\shorthandon{"}}
  }
}
\usetikzlibrary{quotes}
\usepackage{pgfplots}
\usepgfplotslibrary{colormaps}
\pgfplotsset{compat=1.10}



%-- Mathe
\usepackage{mathtools} % beinhaltet amsmath
\mathtoolsset{showonlyrefs, centercolon}
\newtagform{brackets}[\textbf]{[}{]}
\usetagform{brackets}
\usepackage{wasysym}
\usepackage{amssymb} %zusätzliche Symbole
\usepackage{latexsym} %zusätzliche Symbole
\usepackage{stmaryrd} %für Blitz
\usepackage{nicefrac} %schräge Brüche
\usepackage{cancel} %Befehle zum Durchstreichen
\usepackage{mathdots}
% \usepackage{bm}
\DeclareSymbolFont{bbold}{U}{bbold}{m}{n}
\DeclareSymbolFontAlphabet{\mathbbold}{bbold}
\newcommand{\mathds}[1]{\mathbb{#1}} %Um Kompatibilitaet mit frueheren Benutzung von dsfont herzustellen
\newcommand{\ind}{\mathbbold{1}} % charakteristische-Funktion-Eins
\def\mathul#1#2{\color{#1}\underline{{\color{black}#2}}\color{black}} %farbiges Untersteichen im Mathe-Modus

%-- Alles was mit Schrift und XeteX zu tun hat
\ifxetexorluatex
    % XeLaTeX or LuaTeX
	\ifxetex
		\usepackage{mathspec} %beinhaltet fontspec 
	\else
		\usepackage[no-math]{fontspec}
	\fi
    \usepackage{polyglossia} %babel-ersatz
    \setmainlanguage[spelling=new,babelshorthands=false]{german}
	\newcommand\glqq{"}
	\newcommand\grqq{"}
	\defaultfontfeatures{Mapping=tex-text, WordSpace={1.4}} %
    \setmainfont[Ligatures=Common, BoldFont={* Bold}, ItalicFont={* Light Italic},SmallCapsFont={Linux Libertine O}, SmallCapsFeatures={Letters=SmallCaps}]{Source Sans Pro}
	\setsansfont[Scale=MatchLowercase,Ligatures=Common, BoldFont={* Medium}]{Ubuntu}
	\ifxetex
		\setallmonofonts[Scale=MatchLowercase, ItalicFont={*}]{Consolas} 
	\else
		\setmonofont[Scale=MatchLowercase, ItalicFont={*}]{Consolas}
	\fi
	\usepackage{xltxtra}
	\usepackage{fontawesome}
	\usepackage{microtype}
\else
    % default: pdfLaTeX
    \usepackage[ngerman]{babel}
    \usepackage[T1]{fontenc}
    \usepackage[utf8]{inputenc}
	\usepackage{textcomp} %verhindert ein paar Fehler bei den Fonts
    \usepackage[babel=true, tracking=true]{microtype}
	\usepackage{lmodern}
	\renewcommand{\familydefault}{\sfdefault} 
\fi


%--Mixed
\usepackage[neverdecrease]{paralist}
\usepackage[german=quotes]{csquotes}
\usepackage{makeidx}
\usepackage{booktabs}
\usepackage{wrapfig}
\usepackage{float}
\usepackage[margin=10pt, font=small, labelfont={sf, bf}, format=plain, indention=1em]{caption}
\captionsetup[wrapfigure]{name=Abb. }
\usepackage{stackrel}
\usepackage{ifthen}
\usepackage{multicol}
\usepackage{qrcode}
\flushbottom

%--Hyphenation
\hyphenation{di-men-si-o-na-le}


%--Unterstreichung
\usepackage[normalem]{ulem}
\setlength{\ULdepth}{1.8pt}

%--Indexverarbeitung
\newcommand{\bet}[1]{\textbf{#1}}
\newcommand{\Index}[1]{\textbf{#1}\index{#1}}
\makeindex
\setindexpreamble{{\noindent \itshape Die \emph{Seitenzahlen} sind mit Hyperlinks zu den entsprechenden Seiten versehen, also anklickbar} \ifxetexorluatex \faHandUp\fi \par \bigskip}
\renewcommand{\indexpagestyle}{scrheadings}


%--Marginnote & todonotes
\usepackage{marginnote}
\renewcommand*{\marginfont}{\color{dark_gray} \itshape \footnotesize }
\usepackage[textsize=small]{todonotes}

%--Todonotes mit tikz/externalize kompatibel machen
\LetLtxMacro{\oldtodo}{\todo}
\renewcommand{\todo}[2][]{\tikzexternaldisable\oldtodo[#1]{#2}\tikzexternalenable}
\LetLtxMacro{\oldmissingfigure}{\missingfigure}
\renewcommand{\missingfigure}[2][]{\tikzexternaldisable\oldmissingfigure[{#1}]{#2}\tikzexternalenable}

%--Konfiguration von Hyperref pdfstartview=FitH, 
\usepackage[hidelinks, pdfpagelabels,  bookmarksopen=true, bookmarksnumbered=true, linkcolor=black, urlcolor=SkyBlue2, plainpages=false,pagebackref, citecolor=black, hypertexnames=true, pdfauthor={Jannes Bantje}, pdfborderstyle={/S/U}, linkbordercolor=SkyBlue2, colorlinks=false]{hyperref}

% \let\orighref\href
\ifxetexorluatex
\newcommand{\hrefsym}[2]{\href{#1}{\texttt{#2\,\faExternalLink}}}
\newcommand{\hrefsymmail}[2]{\href{#1}{\texttt{\faEnvelope\,#2}}}
\renewcommand{\url}[1]{\hrefsym{#1}{\nolinkurl{#1}}}
\fi
\providecommand{\hrefsym}{\href}
% \usepackage{cleveref}
% \crefformat{footnote}{#2\footnotemark[#1]#3}

%--Römische Zahlen
\newcommand{\RM}[1]{\MakeUppercase{\romannumeral #1{}}}

%--Punkte (nach hyperref)
\usepackage{ellipsis}

%%--Abkürzungen etc.
\newcommand{\light}{\text{\Large $\lightning$}}

%-- Definitionen von weiteren Mathe-Befehlen
\DeclareMathOperator{\re}{Re} %Realteil
\DeclareMathOperator{\im}{Im} %Imaginaerteil
\DeclareMathOperator{\id}{id} %identische Abbildung
\DeclareMathOperator{\Sp}{Sp} %Spur
\DeclareMathOperator{\sgn}{sgn} %Signum
\DeclareMathOperator{\alt}{Alt} %Alternierende n-linearFormen
\DeclareMathOperator{\End}{End} %Endomorphismen
\DeclareMathOperator{\Vol}{Vol} %Volumen
\DeclareMathOperator{\dom}{dom} %Domain
\DeclareMathOperator{\Hom}{Hom} %Homomorphismen
\DeclareMathOperator{\bild}{Bild} %Bild
\DeclareMathOperator{\Kern}{Kern }
\DeclareMathOperator{\rg}{rg} %Rang
\DeclareMathOperator{\Rg}{Rg} %Rang
\DeclareMathOperator{\diam}{diam} %Durchmesser
\DeclareMathOperator{\dist}{dist} %Distanz
\DeclareMathOperator{\grad}{grad} %Gradient
\DeclareMathOperator{\dive}{div} %Gradient
\DeclareMathOperator{\rot}{rot} %Rotation
\DeclareMathOperator{\hess}{Hess} %Hesse-Matrix
\DeclareMathOperator{\koker}{Koker} %Kokern
\DeclareMathOperator{\aut}{Aut} %Automorphismen
\DeclareMathOperator{\ord}{ord} %Ordnung
\DeclareMathOperator{\ggT}{ggT} %ggT
\DeclareMathOperator{\kgV}{kgV} %kgV
\DeclareMathOperator{\Gr}{Gr} %Gerade
\DeclareMathOperator{\Kr}{Kr} %Kreis
\DeclareMathOperator{\Char}{char} %Charakteristik
\DeclareMathOperator{\Aut}{Aut} %Automorphismen
% \DeclareMathOperator{\D}{D} %formale Ableitung
\newcommand{\D}{\ensuremath{\mathrm{D}\mkern-1.5mu}}
\newcommand{\mathd}{\ensuremath{\mathrm{d}\mkern-1.0mu}}
\newcommand{\Tmap}{\ensuremath{\mathrm{T}\mkern-1.0mu}}
\newcommand{\opL}{\ensuremath{\mathrm{L}\mkern-0.6mu}}
\DeclareMathOperator{\cov}{cov} %Kovarianz
\DeclareMathOperator{\Gal}{Gal} %Galois
\DeclareMathOperator{\supp}{supp} %Träger
\DeclareMathOperator{\Sym}{Sym} %Symmetrische Gruppe
\DeclareMathOperator{\Zyl}{Zyl} %Zylinder
\DeclareMathOperator{\Mod}{Mod} %Moduln
\DeclareMathOperator{\EPK}{EPK} %Einpunktkompaktifizierung
\DeclareMathOperator{\conj}{conj} %Konjugation
\DeclareMathOperator{\ann}{ann} %Annulator
\DeclareMathOperator{\tor}{tor} %Torsionsmodul
\DeclareMathOperator{\Int}{Int} %Interior/Inneres
\DeclareMathOperator{\Br}{Br} %Brauerergruppe
\DeclareMathOperator{\GL}{GL} %allgemeine lineare Gruppe
\DeclareMathOperator{\SL}{SL} %spezielle lineare Gruppe
\DeclareMathOperator{\SU}{SU} %spezielle unitäre Gruppe
\DeclareMathOperator{\Tr}{Tr} %Spur/Trace
\DeclareMathOperator{\conv}{conv} %Konvexe Teilmenge
% \DeclareMathOperator{\tr}{tr}
\newcommand{\tr}{\mathrm{tr}}
\newcommand{\w}{\mathrm{w}}
\newcommand{\weakT}[1]{\ensuremath{\mathcal{T}_{#1}^{\mkern+1.0mu\text{\raisebox{0.4ex}{$\mathrm{w}$}}}}}
\newcommand{\weakTstar}[1]{\ensuremath{\mathcal{T}_{#1}^{\mkern+1.0mu\text{\raisebox{0.4ex}{$\mathrm{w}$}}^*}}}
\newcommand{\rand}[1]{\ensuremath{\partial^{\scriptscriptstyle #1}}}


% -- Kategorien
\DeclareMathOperator{\Ob}{Ob} %Objekte
\DeclareMathOperator{\Mor}{Mor} %Morphismen

% \DeclareMathOperator{\TOP}{TOP} %Topologische Räume
\newcommand{\TOP}{\textsc{Top}}
% \DeclareMathOperator{\HTOP}{HTOP} %Topologische Räume + Homotopieäquivalenz
\newcommand{\HTOP}{\textsc{HTop}}
% \DeclareMathOperator{\VR}{VR} %Vektorräume
\newcommand{\VR}{\textsc{VR}}
% \DeclareMathOperator{\MOD}{MOD} %Moduln
\newcommand{\MOD}{\textsc{Mod}}
% \DeclareMathOperator{\MENGEN}{MENGEN} %Mengen
\newcommand{\MENGEN}{\textsc{Mengen}}
% \DeclareMathOperator{\MAN}{MAN} %Mannigfaltigkeiten
\newcommand{\MAN}{\textsc{Man}}
% \DeclareMathOperator{\GRUPPEN}{GRUPPEN} %Gruppen
\newcommand{\GRUPPEN}{\textsc{Gruppen}}
% \DeclareMathOperator{\ABELGRUPPEN}{ABEL.GRUPPEN} %abelsche Gruppen
\newcommand{\ABELGRUPPEN}{\textsc{Abel.Gruppen}}
% \DeclareMathOperator{\KAT}{KAT} %Kategorien
\newcommand{\KAT}{\textsc{Kat}}
% \DeclareMathOperator{\FUN}{FUN} %Funktoren
\newcommand{\FUN}{\textsc{Fun}}


\newcommand{\op}{\mathrm{op}}

\newcommand{\Ker}{\ensuremath{\Kern \,}}
\newcommand{\zirkel}{\rotatebox[origin = c]{-90}{$\varangle$}}
\newcommand{\wirkung}{\!\curvearrowright\!}

\newcommand{\oE}{\OE}

%--Beweisende
\newcommand{\bewende}{\ifmmode \tag*{$\square$} \else \hfill $\square$ \fi}

%--Volumen
\newcommand{\vol}[1]{\ensuremath{\Vol_{#1}}}

%--Skalarprodukt
\DeclarePairedDelimiterX\sprod[2]{\langle}{\rangle}{#1\,\delimsize\vert\,#2}

%--Betrag, Gaußklammer
\DeclarePairedDelimiter{\abs}{\lvert}{\rvert}
\DeclarePairedDelimiter{\ceil}{\lceil}{\rceil}

%--Norm
\DeclarePairedDelimiter\doppelstrich{\Vert}{\Vert}
\newcommand{\norm}[2][\relax]{
\ifx#1\relax \ensuremath{\doppelstrich*{#2}}
\else \ensuremath{\doppelstrich*{#2}_{#1}}
\fi}


%--Umklammern
\DeclarePairedDelimiter\enbrace{(}{)}
\DeclarePairedDelimiter\benbrace{[}{]}
\DeclarePairedDelimiter\lenbrace{<}{>}

%--Mengen
\DeclarePairedDelimiterX\mengenA[1]{\lbrace}{\rbrace}{#1}
\DeclarePairedDelimiterX\mengenB[2]{\lbrace}{\rbrace}{#1\, \delimsize\vert \, #2}

\makeatletter
\newcommand{\set}[2][\relax]{
\ifx#1\relax \ensuremath{
\mengenA*{#2}}
\else \ensuremath{%
  \mengenB*{#1}{#2}}
\fi}
\makeatother


%--offen und abgeschlossen
\newcommand{\off}{\! \stackrel[\text{offen}]{}{\subset} \!}
\newcommand{\abg}{\! \stackrel[\text{abg.}]{}{\subset} \!}

%--Differential
\newcommand{\diff}[2]{\ensuremath{\frac{{\partial #1}}{{\partial #2}} }}
\newcommand{\diffd}[2]{\ensuremath{\frac{\mathd #1}{\mathd #2} }}

%--Diagonale Linien für Matrizen
\newcommand\tikzmark[1]{%
  \tikzexternaldisable\tikz[overlay,remember picture,baseline] \node [anchor=base] (#1) {};\tikzexternalenable}

\newcommand\MyLine[3][]{%
  \tikzexternaldisable\begin{tikzpicture}[overlay,remember picture]
    \draw[#1] (#2.south) -- (#3.north);
  \end{tikzpicture}\tikzexternalenable}
  
%--Jordankasten
\newcommand{\Jord}{ \tikzexternaldisable \begin{BMAT}{|ccc|}{|ccc|} %
				0 & & \\ %
				1\tikzmark{z} &  \diagdown & \\ %
				& \tikzmark{u}1 \MyLine[thick]{z}{u} &  0 %
			\end{BMAT} \tikzexternalenable}
			
%--Jordankasten mit Lambda
\newcommand{\JordLa}[1][\empty]{%
	\tikzexternaldisable\ifthenelse{\equal{#1}{\empty}}
		{\begin{BMAT}{|ccc|}{|ccc|} %
				\lambda & & \\ %
				1\tikzmark{z} &  \diagdown & \\ %
				& \tikzmark{u}1 \MyLine[thick]{z}{u} &  \lambda  %
			\end{BMAT}}
		{\begin{BMAT}{|ccc|}{|ccc|} %
				\lambda_{#1}  & & \\ %
				1\tikzmark{z} &  \diagdown & \\ %
				& \tikzmark{u}1 \MyLine[thick]{z}{u} &  \lambda_{#1}  %
			\end{BMAT}}\tikzexternalenable}

%--Inhaltsverzeichnis
\usepackage[tocindentauto]{tocstyle}
\usetocstyle{KOMAlike}			

\newcommand{\fach}{$K$-Theorie und die Hopf-Invariante}
% \newcommand{\shortFach}{Analysis, Topologie, Geometrie}
\newcommand{\semester}{SoSe 2015}
\newcommand{\homepage}{https://wwwmath.uni-muenster.de/reine/u/topos/lehre/SS2015/KTheorie-Hopf/Hopf.html}

\newcommand{\prof}{Dr.\ Ulrich Penning}
\providecommand{\verfasser}{Jannes Bantje}

%--Konfiguration von scrheadings
\setheadsepline{1pt}[\color{light_gray}]
\pagestyle{scrheadings}
%\fancyhf[H, F]{}
\clearscrheadfoot

\providecommand{\shortFach}{\fach}
\lehead{ \includegraphics[height=0.5 cm,keepaspectratio]{../!config/Bilder/Logo_WWU_Muenster_light_gray.pdf}}
\rehead{\rule{0cm}{0.5cm}\footnotesize \sffamily \color{light_gray} \verfasser -- Mitschrift \shortFach}
\lohead{\rule{0cm}{0.5cm} \footnotesize \sffamily \color{light_gray} Stand: \today \; \thistime[:]}
\rohead{\includegraphics[height=0.5 cm,keepaspectratio]{../!config/Bilder/fb10logo_gray.pdf}}


\ofoot[{ \color{dark_gray} \LARGE \sffamily \thepage}]{{ \color{dark_gray} \LARGE \sffamily \thepage}} %hier wir auch der plain Stil bearbeitet!
\automark{section}
\ifoot{ \color{dark_gray} \small \leftmark}

%--Metadaten
\author{\verfasser}
\titlehead{\includegraphics[height=1.5cm, keepaspectratio]{../!config/Bilder/Logo_WWU_Muenster.pdf}%
\hfill \includegraphics[height=1.3cm, keepaspectratio]{../!config/Bilder/fb10logo.pdf}}
\title{Skript \fach}
\subtitle{Mitschrift der Vorlesung  \enquote{\fach} von \prof}

\publishers{
Aktuelle Version verfügbar bei: \bigskip\\
\normalsize
\begin{minipage}[t]{0.45\textwidth}
	\centering{
	\qrcode[height=3.5cm, version=4]{https://github.com/JaMeZ-B/latex-wwu}
	\medskip\\
	\includegraphics[height=0.4cm, keepaspectratio]{../!config/Bilder/github_octo.pdf} 
	\includegraphics[height=0.4cm, keepaspectratio]{../!config/Bilder/GitHub_Logo.pdf} (inklusive Sourcecode)\\
	{\small\url{https://github.com/JaMeZ-B/latex-wwu}}}
\end{minipage}
\hfill
\begin{minipage}[t]{0.45\textwidth}
	\centering{
	\qrcode[height=3.5cm, version=4]{btsync://B6WH2DISQ5QVYIRYIEZSF4ZR2IDVKPN3I?n=latex_share}
	\medskip\\
	\includegraphics[height=0.4cm, keepaspectratio]{../!config/Bilder/bt_sync_logo.pdf} 
	{\large \textbf{Bittorrent} Sync} \\
	\texttt{B6WH2DISQ5QVYIRYIEZSF4ZR2IDVKPN3I}
	}
\end{minipage}
}


\begin{document}
\pagenumbering{Roman}
\maketitle
\begin{abstract}
\section*{Vorwort --- Mitarbeit am Skript}
Dieses Dokument ist eine Mitschrift aus der Vorlesung \enquote{\fach, \semester}, gelesen von \prof. Der Inhalt entspricht weitestgehend dem Tafelanschrieb. Für die
Korrektheit des Inhalts übernehme ich keinerlei Garantie! Für Bemerkungen und Korrekturen -- und seien es nur Rechtschreibfehler -- bin ich sehr dankbar. 
Korrekturen lassen sich prinzipiell auf drei Wegen einreichen: 
\begin{itemize}
	\item Persönliches Ansprechen in der Uni, Mails an \hrefsymmail{mailto:\mail}{\mail} (gerne auch mit annotieren PDFs) 
	\item \emph{Direktes} Mitarbeiten am Skript: Den Quellcode poste ich auf GitHub (siehe oben), also stehen vielfältige Möglichkeiten der Zusammenarbeit zur Verfügung:
	Zum Beispiel durch Kommentare am Code über die Website und die Kombination Fork + Pull Request. Wer sich verdient macht oder ein Skript zu einer Vorlesung, die 
	ich nicht besuche, beisteuern will, dem gewähre ich gerne auch Schreibzugriff.
	
	Beachten sollte man dabei, dass dazu ein Account bei \url{github.com} notwendig ist, der allerdings ohne Angabe von persönlichen Daten angelegt werden kann. 
	Wer bei GitHub (bzw. dem zugrunde liegenden Open-Source-Programm \enquote{\texttt{git}}) -- verständlicherweise -- Hilfe beim Einstieg braucht, dem helfe ich gerne 
	weiter. Es gibt aber auch zahlreiche empfehlenswerte Tutorials im Internet.\footnote{zB. \url{https://try.github.io/levels/1/challenges/1}, ist auf Englisch, aber dafür 
	interaktives LearningByDoing}
	\item \emph{Indirektes} Mitarbeiten: \TeX-Dateien per Mail verschicken. 
	
	Dies ist nur dann sinnvoll, wenn man einen ganzen Abschnitt ändern möchte (zB. einen alternativen Beweis geben), da ich die Änderungen dann per Hand einbauen muss! Ich freue mich aber auch über solche Beiträge!
\end{itemize}
\section*{Vorlesungshomepage}
{\centering 
\begin{minipage}[c][][c]{\textwidth}
	\centering \qrcode[height=3cm]{\homepage} \medskip\\
	\footnotesize \url{\homepage}
\end{minipage}
\par}

\end{abstract}

\tableofcontents
\cleardoubleoddemptypage

\pagenumbering{arabic}
\setcounter{page}{1}

\setcounter{section}{-1}
\section{Einführung} % (fold)
\label{sec:0}

\subsection[Definition: Reelle Divisionsalgebra]{Definition} % (fold)
\label{sub:01}
Eine $\mathds{R}$-Algebra $\mathcal{A}$ mit $1$, die nicht notwendigerweise assoziativ ist, heißt \Index{reelle Divisionsalgebra}, falls jedes Element $a\not= 0$ invertierbar 
ist.
\minisec{Beispiele}
Die reellen Zahlen $\mathds{R}$, die komplexen Zahlen $\mathds{C}$, die Quaternionen\footnote{siehe auch \url{https://de.wikipedia.org/wiki/Quaternion}} $\mathds{H}$ und die sogenannten Cayley-Zahlen\footnote{auch Oktonionen oder reelle Okataven; siehe  \url{http://de.wikipedia.org/wiki/Oktave_(Mathematik)}} $\mathds{O}$. Dabei ist
$\mathds{O} \cong \mathds{H} \oplus \mathds{H}$, wobei die Multiplikation wie folgt definiert ist:
\[
	(a,b) \cdot (d,c) = (ac - d^* b, da + bc^*) 
\]
Dabei ist $d= x + iy + jz + kw$ und $d^* = x -iy -jz - kw$.
% subsection 11 (end)

\subsection{Theorem 1.1 (Adams)} % (fold)
\label{sub:02}
Folgende Aussagen sind äquivalent:
\begin{enumerate}[(i)]
	\item $\mathds{R}^n$ besitzt die Struktur einer reellen Divisionsalgebra.
	\item Entweder ist $n=1$ oder $n \ge 2$ ist gerade und es gibt eine stetige Abbildung $f \colon S^{2n-1} \to S^n$ mit der sogenannten \Index{Hopf-Invariante} $1$.
	\item Es gilt $n \in \set{1,2,4,8}$.
\end{enumerate}
% subsection 11 (end)

\subsection{Geplanter Verlauf der Vorlesung} % (fold)
\label{sub:inhalt}
Wir werden diesen Satz mit Hilfe der sogenannten $K$-Theorie beweisen. Dazu müssen wir aber einiges an Vorarbeit leisten:
\begin{enumerate}[1.]
	\item Verallgemeinerte Kohomologietheorien
	\item Vektorbündel
	\item $K$-Theorie
	\begin{enumerate}[label*=\arabic*.]
		\item Produkte in $K$-Theorie
		\item Bott-Periodizität
		\item Thom-Isomorphismus
	\end{enumerate}
	\item Hopf-Invariante und der Beweis des Satzes
\end{enumerate}
% subsection inhalt (end)
% section 0 (end)
\newpage

\section{Kohomologietheorien} % (fold)
\label{sec:1}

\subsection[Definition: Kontravarianter Funktor]{Definition} % (fold)
\label{sub:11}
Seien $\mathcal{C}$ und $\mathcal{D}$ Kategorien. Ein \Index{kontravarianter Funktor} $F \colon \mathcal{C} \to \mathcal{D}$ 
\begin{itemize}
	\item ordnet jedem Objekt $c \in \mathrm{Obj}(\mathcal{C})$ ein Objekt $F(c) \in \mathrm{Obj}(\mathcal{D})$ zu
	\item ordnet jedem Morphismus $f \colon c \to c'$ in $\mathrm{Mor}_{\mathcal{C}}(c,c')$ einen Morphismus $F(f) \colon F(c') \to F(c)$ zu 
\end{itemize}
sodass folgende Eigenschaften gelten:
\begin{itemize}
	\item Für $f \colon c \to c'$ und $g \colon c' \to c''$ gilt $F(g \circ f) = F(f) \circ F(g)$
	\item $F(\id_c) = \id_{F(c)}$
\end{itemize}
\minisec{Beispiel (Dualraumfunktor)}
Sei $K$ ein Körper. Sei $\VEKT_K$ die Kategorie der endlich dimensionalen $K$-Vektorräume und linearen Abbildungen. Sei 
${ }^* \colon \VEKT_K \to \VEKT_K$ gegeben durch $V \mapsto \Hom_K(V,K)=V^*$. Für $f \colon V \to W$ sei $f^* \colon W^* \to V*$, 
$\varphi \mapsto \varphi \circ f$. Dies ist ein kontravarianter Funktor. Insbesondere ist $(f \circ g)^* = g^* \circ f^*$.
% subsection 11 (end)

\subsection[Definition: Verallgemeinerte Kohomologietheorie]{Definition} % (fold)
\label{sub:12}
Sei $\TOP^2$ die Kategorie der Raumpaare $(X,A)$. Sei $V \colon \TOP^2 \to \TOP^2$ der Funktor $(X,A) \mapsto(A,\emptyset)$. Sei $R$ ein kommutativer Ring mit $1$. Eine 
\Index{verallgemeinerte Kohomologietheorie} $h^* = (h^n, \partial^n)_{n \in \mathds{Z}}$ mit Werten in $R$-Moduln ist ein kontravarianter Funktor 
$h^* \colon \TOP^2 \to \textsc{Gr}\text{-}R\text{-}\MOD$ zusammen mit einer natürlichen Transformation 
\[
	\partial^* \colon h^* \circ V \to h^{*+1},
\]
so dass die folgenden Eigenschaften gelten:
\begin{enumerate}[a)]
	\item \Index{Homotopieinvarianz}: Seien $f,g \colon (X,A) \to (Y,B)$ homotope Abbildungen von Raumpaaren. Dann gilt für alle $n \in \mathds{Z}$: 
	\[
		h^n(f)= h^n(g)
	\]
	\item \Index{Lange exakte Paarsequenz}: Sei $(X,A)$ ein Raumpaar. Seien $(A,\emptyset) \xhookrightarrow{i}(X,\emptyset)$ und $j \colon (X,\emptyset) \to (X,A)$ die 
	kanonischen Inklusionen. Dann ist die Folge 
	\[
		\begin{tikzcd}
			\cdots \rar["\partial^{n-1}"] & h^n(X,A;R) \rar["h^{n}(j)"] & h^n (X,\emptyset;R) \rar["h^n(i)"] & h^n(A, \emptyset;R) \rar["\partial^n"] & h^{n+1}(X,A;R) \rar & 
			\cdots
		\end{tikzcd}
	\]
	\item \Index{Ausschneidung}: Sei $(X,A)$ ein Raumpaar, $U \subset A$, sodass $\overline{U} \subseteq \mathring{A}$. Dann ist die von der Inklusion 
	$\iota \colon (X \setminus U, A \setminus U) \hookrightarrow (X,A)$ induzierte Abbildung $h^n(\iota) \colon h^n(X,A;R) \to h^n(X \setminus U, A \setminus U;R)$ ein Isomorphismus.
\end{enumerate}
\minisec{Bemerkungen}
\begin{enumerate}[(i)]
	\item Wir schreiben kurz $h^n(X,A)$ für $h^n(X,A;\mathds{Z})$ und $f^*$ für $h^n(f)$, falls sich der Grad aus dem Kontext ergibt. 
	\item $h^n(\set{\mathrm{pt}}) =: h^n$ heißen \Index{Koeffizienten der Kohomologietheorie}.
	\item Einige Kohomologietheorien besitzen ein \Index{externes Produkt}, das heißt eine natürliche Transformation\todo{Diagramm hinzufügen}
	\[
		h^n(X,A;R) \times h^m(Y,B;R) \longrightarrow h^{n+m} \enbrace[\big]{(X,A) \times (Y,B);R} 
	\]
	die bilinear und assoziativ ist. 
	Dabei ist $(X,A) \times (Y,B) = \enbrace[\big]{X \times Y, (A \times Y) \cup (X \times B)}$.
	
	Falls ein externes Produkt für $h^*$ existiert, dann ist $h^*(X,A;R) = \bigoplus_{n \in \mathds{Z}} h^n(X,A;R)$ ein \Index{graduierter Ring} mit der Multiplikation 
	\[
		\begin{tikzcd}[row sep=large, column sep=small]
			h^n(X,A;R) \times h^m(X,A;R) \arrow[dashed]{rr} \arrow["\times"']{dr} & & h^{n+m}(X,A;R) \\
			& h^{n+m} \enbrace[\big]{(X,A) \times (X,A);R} \arrow["h^{n+m}(\Delta)"']{ur} &
		\end{tikzcd}
	\]
\end{enumerate}
% subsection 12 (end)

\subsection[Definition: Ko-Raumpaar und Kofaserung]{Definition} % (fold)
\label{sub:13}
Ein Raumpaar $(X,A)$ heißt \Index{Ko-Raumpaar}, falls 
\begin{itemize}
	\item $X$ ein kompakter Hausdorffraum ist
	\item $A \subset X$ abgeschlossen ist
	\item die Inklusion $A \hookrightarrow X$ eine \Index{Kofaserung} ist, das heißt für jede stetige Abbildung\marginnote{Homotopieerweiterungseigenschaft} 
	$f \colon A \times I \cup X \times \set{0} \to Y$ existiert 
	$H \colon X \times I \to Y$, sodass das folgende Diagramm kommutiert:
	\[
		\begin{tikzcd}[sep=large]
			A \times I \cup X \times \set{0} \rar["f"] \dar[hook]& Y \\
			X \times I \urar["H"'] & 
		\end{tikzcd}
	\] 
\end{itemize}
Mit $\TOP^2_{\mathrm{Kof}} \subseteq \TOP^2$ bezeichnen wir die Kategorie der Ko-Raumpaare und stetigen Abbildungen.
\minisec{Beispiel}
\begin{itemize}
	\item Sei $X$ ein kompakter \CW-Komplex, $A \subset X$ ein Unterkomplex. Dann ist $\iota \colon A \hookrightarrow X$ eine Kofaserung.
	\item $X$ eine kompakte Mannigfaltigkeit, $A$ eine abgeschlossene Untermannigfaltigkeit.
\end{itemize}
\minisec{Bemerkung}
Falls $(X, \set{x_0})$ ein Objekt in $\TOP^2_{\mathrm{Kof}}$ ist, dann heißt $X$ \Index{wohlpunktiert}.
\minisec{Beispiel}
$X= [-1,1] \times [-1,1]$, $x_0=(0,0)$, $Y=X$. Sei $\gamma \colon I \to X$ ein Pfad mit $\gamma(0)= (0,0)=x_0$. Wir erhalten eine Abbildung 
$f \colon \set{x_0} \times I \cup  X \times \set{0} \to X $
\missingfigure{Skizze zu der Abbildung}
Kofaserung heißt in diesem Fall, dass der Raum $X$ entlang des Pfades \enquote{mitgezogen} werden kann.
% subsection 13 (end)
% section 1 (end)
\newpage

\section{Vektorbündel} % (fold)
\label{sec:2}
Sei für den Verlauf dieses Kapitels $K \in \set{\mathds{R},\mathds{C}}$. 

\subsection[Definition: $K$-Vektorbündel]{Definition} % (fold)
\label{sub:21}
Sei $X$ ein topologischer Raum. Eine stetige surjektive Abbildung $p \colon E \to X$ heißt $K$-\Index{Vektorbündel}, falls gilt: 
\begin{enumerate}[a)]
	\item $p^{-1}(\set{x}) =: E_x$ ist ein $K$-Vektorraum für alle $x \in X$. $E_x$ heißt die \Index{Faser} von $E$ über $x$. 
	\item Für alle $x \in X$ gibt es eine Umgebung $U$, ein $n \in \mathds{N}$ und einen Homöomorphismus $\varphi_U \colon p^{-1}(U) \to U \times K^n$, sodass das folgende 
	Diagramm kommutiert
	\[
		\begin{tikzcd}
			p^{-1}(U) \arrow["\varphi_U"]{rr} \drar["p"'] & & U \times K^n  \dlar["\mathrm{pr}_U"] \\
			& U 
		\end{tikzcd}
	\]
	und die Einschränkung $\varphi_U\big|_{E_x} \colon E_x \to K^n$ für alle $x \in U$ ein Isomorphismus von $K$-Vektorräumen ist. Diese Eigenschaft wird oft als lokale 
	Trivialität bezeichnet.
\end{enumerate}
\minisec{Beispiel}
\begin{itemize}
	\item Sei $X$ ein beliebiger topologischer Raum. Dann ist $X \times K^n \xrightarrow{\mathrm{pr}_X} X$ das triviale $K$-Vektorbündel über $X$.
	\item Betrachte $S^n \subset \mathds{R}^{n+1}$. Sei $\Tmap S^n = \set[(x,v) \in S^n \times \mathds{R}^{n+1}]{\skal*{x}{v}=0}$. Definiere $p \colon \Tmap S^n \to S^n$, 
	$(x,v) \mapsto x$. Damit erhalten wir ein $\mathds{R}$-Vektorbündel über $S^n$. 
	
	Übung: Zeige lokale Trivialität.
\end{itemize}
% subsection 21 (end)

\subsection[Definition: Vektorraumbündelmorphismus]{Definition} % (fold)
\label{sub:22}
Seien $p \colon E \to X$ und $p' \colon E' \to X'$ $K$-Vektorbündel. Ein \bet{Morphismus von Vektorbündeln}\index{Morphismus!von Vektorbündeln} oder 
\Index{Vektorraumbündelmorphismus} von $E$ nach $E'$ ist ein Paar $(f, \overline{f})$ von stetigen Abbildungen $\overline{f} \colon E \to E'$ und $f \colon X \to X'$, sodass 
das folgende Diagramm kommutiert
\[
	\begin{tikzcd}
		E \rar["\overline{f}"] \dar["p"'] & E' \dar["p'"] \\
		X \rar["f"] & X'
	\end{tikzcd}
\]
und sodass die Einschränkung $\overline{f}\big|_{E_x} \colon E_x \to E'_{f(x)}$ für jedes $x \in X$ ein $K$-Vektorraumhomomorphismus ist. Wir bezeichnen die Kategorie der 
endlich dimensionalen $K$-Vektorraumbündel über einem topologischen Raum $X$ zusammen mit den Morphismen $(\id_X, \overline{f})$ mit $\VEKT_K(X)$.
% subsection 22 (end)
% section 2 (end)

\cleardoubleoddemptypage
\pagenumbering{Alph}
\setcounter{page}{1}
\printindex
\listoffigures
\todototoc
\listoftodos[To-do's und andere Baustellen]
\end{document}